\documentclass[aspectratio=169,10pt]{beamer}
\usetheme[numbering=fraction,progressbar=frametitle]{metropolis}
\usepackage{fontspec}
\setsansfont{Latin Modern Sans}
\usepackage{amsmath,amssymb,mathtools,graphicx,booktabs}
\usepackage{polyglossia}
\setdefaultlanguage{french}
\setbeamertemplate{navigation symbols}{}
\setbeamerfont{frametitle}{size=\large}
\setbeamerfont{normal text}{size=\normalsize}
\setlength{\parskip}{.45em}
\hypersetup{pdfauthor={},pdfsubject={Lecture slides},pdfkeywords={},pdftitle={Dérivation : du taux à la tangente}}
\providecommand{\N}{\mathbb{N}}
\providecommand{\Z}{\mathbb{Z}}
\providecommand{\D}{\mathbb{D}}
\providecommand{\Q}{\mathbb{Q}}
\providecommand{\R}{\mathbb{R}}
\providecommand{\C}{\mathbb{C}}
\providecommand{\vect}[1]{\overrightarrow{#1}}
\providecommand{\vi}{\vec{\imath}}
\providecommand{\vj}{\vec{\jmath}}
\providecommand{\vk}{\vec{k}}
\providecommand{\norm}[1]{\left\lVert #1\right\rVert}
\providecommand{\vc}[2]{\begin{pmatrix}#1\\#2\end{pmatrix}}
\providecommand{\deter}[2]{\det\left(#1,#2\right)}
\providecommand{\abs}[1]{\left\lvert #1\right\rvert}
\providecommand{\dd}{\,\mathrm{d}}
\providecommand{\e}{\mathrm{e}}
\providecommand{\Card}{\operatorname{Card}}
\providecommand{\E}{\operatorname{E}}
\providecommand{\V}{\operatorname{V}}
\providecommand{\degre}{^{\circ}}
\providecommand{\un}[1]{\,\mathrm{#1}}
\providecommand{\sci}[2]{#1\times10^{#2}}
\providecommand{\pgcd}{\operatorname{PGCD}}
\providecommand{\ppcm}{\operatorname{PPCM}}
\title{Dérivation : du taux à la tangente}
\subtitle{Cours et exemples}
\date{}
\author{}
\begin{document}
\begin{frame}[plain]
\vfill
{\usebeamerfont{title}\inserttitle\par}
\medskip
{\usebeamerfont{subtitle}\insertsubtitle\par}
\medskip
{\usebeamercolor[fg]{progress bar}\rule{\textwidth}{.5pt}}
\vfill
\end{frame}
\begin{frame}{Le parcours du chapitre}
\textbf{Objectifs.} Interpréter un taux ; approcher un nombre dérivé ; déterminer une tangente ; comprendre les échecs de dérivabilité.

\textbf{Prérequis.} Images, pente d’une droite, calcul littéral, racine carrée et valeur absolue.

\textbf{Repère annuel.} 3-12 novembre : 8 h
\end{frame}
\begin{frame}{Taux de variation et sécantes}
Pour \(a\ne b\) dans le domaine de f, le taux de variation entre a et b vaut \(\dfrac{f(b)-f(a)}{b-a}\). C’est la pente de la sécante passant par les points correspondants. En contexte, l’unité du taux est l’unité de f divisée par celle de la variable.
\end{frame}
\begin{frame}{Taux de variation et sécantes}
Pour \(f(x)=x^2\), entre 1 et 3, le taux vaut \(\dfrac{9-1}{3-1}=4\). Entre 1 et \(1+h\), il vaut \(2+h\). Les pentes pour \(h=0{,}1\) et \(h=-0{,}1\) sont 2,1 et 1,9.
\end{frame}
\begin{frame}{Taux de variation et sécantes}
Le taux n’est pas \(\dfrac{f(b)}{f(a)}\) ; h représente un écart, positif ou négatif selon le domaine.
\end{frame}
\begin{frame}{Taux de variation et sécantes}
Calculer le taux de \(f(x)=x^2+3x\) entre 2 et \(2+h\).
\end{frame}
\begin{frame}{Nombre dérivé}
f est dérivable en a lorsque le taux \(\dfrac{[f(a+h)-f(a)]}{h}\) tend vers un réel quand h tend vers 0, dans le domaine. Ce réel est \(f\prime(a)\). On emploie ici une approche intuitive de la limite. En cinématique, si f donne une position en mètres et le temps en secondes, le nombre dérivé représente une vitesse instantanée en \(\mathrm{m\,s^{-1}}\).
\end{frame}
\begin{frame}{Nombre dérivé}
Pour \(s(t)=t^2+3t\), le taux autour de 2 vaut \(7+h\), donc \(s\prime(2)=7\ \mathrm{m\,s^{-1}}\). La vitesse moyenne sur \([2;3]\) vaut \(8\ \mathrm{m\,s^{-1}}\) ; elle n’est pas la vitesse à l’instant 2.
\end{frame}
\begin{frame}{Nombre dérivé}
On ne remplace pas h par 0 dans un quotient non simplifié : on obtiendrait une division par 0.
\end{frame}
\begin{frame}{Nombre dérivé}
Pour \(f(x)=3x^2-2x\), calculer \(f\prime(1)\) à partir du taux.
\end{frame}
\begin{frame}{Tangente et approximation affine}
Si f est dérivable en a, la tangente au point \((a;f(a))\) a pour équation \(y=f(a)+f\prime(a)(x-a)\). Pour h petit, on utilise \(f(a+h)\approx f(a)+f\prime(a)h\). Cette approximation locale ne fournit pas à elle seule une borne de l’erreur.
\end{frame}
\begin{frame}{Tangente et approximation affine}
\begin{columns}[T,onlytextwidth]
\begin{column}{.57\textwidth}
Pour \(f(x)=x^2\) en 2, tangente \(y=4+4(x-2)=4x-4\). On approche \(2{,}01^2\) par \(4+4(0{,}01)=4{,}04\) ; la valeur exacte est \(4{,}0401\), écart \(0{,}0001=h^2\).
\end{column}
\begin{column}{.4\textwidth}
\centering
\includegraphics[width=\linewidth,height=4.5cm,keepaspectratio]{figures/premiere-04-1.pdf}
\end{column}
\end{columns}
\end{frame}
\begin{frame}{Tangente et approximation affine}
Une tangente peut recouper la courbe ailleurs. « Un seul point commun » n’est pas une définition.
\end{frame}
\begin{frame}{Tangente et approximation affine}
On donne \(f(3)=5\) et \(f\prime(3)=-2\). Écrire la tangente et approcher \(f(3{,}04)\).
\end{frame}
\begin{frame}{Quand la dérivabilité échoue}
La fonction racine carrée est définie sur \([0;+\infty[\) mais n’est pas dérivable en 0. La fonction valeur absolue est définie sur \(\R\) mais n’est pas dérivable en 0. Être définie ou avoir une courbe continue ne garantit donc pas la dérivabilité.
\end{frame}
\begin{frame}{Quand la dérivabilité échoue}
Pour \(|x|\), les sécantes issues de l’origine ont toutes pente 1 à droite et \(-1\) à gauche. Le point anguleux ne possède pas une tangente de pente finie unique. Pour la racine, la direction limite est verticale ; il n’y a pas de nombre dérivé réel.
\end{frame}
\begin{frame}{Quand la dérivabilité échoue}
Une pente très grande n’est pas une preuve de non-dérivabilité ; il faut expliquer le comportement du taux.
\end{frame}
\begin{frame}{Quand la dérivabilité échoue}
Justifier que \(g(x)=|x-2|\) n’est pas dérivable en 2. Étudier aussi la dérivabilité de \(x\mapsto\sqrt{x}\) en 0 par son taux.
\end{frame}
\end{document}
